% 6.电路基础.tex 子文件，book文档类，主文档通过\input{6.电路基础.tex}导入
\chapter{第七章 频率特性和谐振现象}

\section{网络函数和频率特性}
若将所关心的网络某一电压或电流作为响应，比例系数记作
\[
\]
称为网络函数。

如果激励和响应属于。

集中参数正弦，其网络函数为
\[
\]

可见。

仍，即
\[
\]
由此可得
\[
\]
式中。

\section{串联谐振电路}
根据谐振的定义
\[
\]
即
\[
\]

由电感或电容分别为
\[
\]
式中$\omega_0$称为$RLC$串联电路的谐振角频率。

发生串联谐振时，令
\[
\]
$\rho$称为$RLC$串联电路的特性阻抗。将谐振角频率，得
\[
\]
又令
\[
\]
称为$RLC$串联电路的品质因数。可得品质因数的另外计算公式
\[
\]
当外加电压$U$不变时，电路中$I、U_L、U_C$与$\omega$的关系是
\[
\]
根据。

从。

当，记为
\[
\]
电流
\[
\]
\[
\]
引用$RLC$串联有效值
\[
\]
可见。

如果$RLC$串联电路中电阻$R$趋于零。

\section{RLC串联电路的频率特性}
当以电阻电压为响应时，
\[
\]
将谐振
\[
\]
\[
\]
由上式。

由幅频特性，需满足
\[
\]
由此解得
\[
\]
\[
\]
截止频率，的关系是
\[
\]
可见，$RLC$带通电路的通带宽度与品质因数成反比。

如果，其转移电压比为
\[
\]
\[
\]
幅频特性和相频特性分别为
\[
\]
当。

从幅频特性。

图。

最后。其转移电压比为
\[
\]
\[
\]
幅频特性和相频特性分别为
\[
\]
\[
\]
当。

图。

\section{并联谐振电路}
由电流源。

GCL并联
\[
\]
实现谐振的条件是。

\endinput